%%%%%%%%%%%%%%%%%%%%%%%%%%%%%%%%%%%%%%%%%
% Beamer Presentation - LaTeX Template
% Version 2.0 (March 8, 2022)
% Original Template: https://www.LaTeXTemplates.com
% Author: Vel (vel@latextemplates.com)
% License: CC BY-NC-SA 4.0
%%%%%%%%%%%%%%%%%%%%%%%%%%%%%%%%%%%%%%%%%

%%%%%%%%%%%%%%%%%%%%%%%%%%%%%%%%%%%%%%%%%
% Este modelo de apresentação foi 
% criado a partir do modelo de Giovanni Spadaro.
% Disponível em: 
% https://github.com/Giovo17/presentation-template-unict-lm-data
%
% Adaptado por Lucas Amaral Taylor para criar uma versão especial 
% para os alunos de Matemática e Estatística da USP (IME-USP).
% Disponível em:
% https://github.com/lucasamtaylor01/IME-template
%%%%%%%%%%%%%%%%%%%%%%%%%%%%%%%%%%%%%%%%%

\documentclass[
	11pt, % Tamanho padrão da fonte
	% t, % Alinhar verticalmente ao topo
	% aspectratio=169, % Definir proporção 16:9
]{beamer}

% Caminho para imagens
\graphicspath{{img/}}

% Pacotes adicionais
\usepackage[alf]{abntex2cite} % Citações ABNT
\usepackage{booktabs} % Linhas de tabela aprimoradas
\usepackage{palatino} % Fonte Palatino
\usepackage[default]{opensans} % Fonte Open Sans

%----------------------------------------------------------------------------------------
%	SELEÇÃO DE LAYOUT E CORES
%----------------------------------------------------------------------------------------

\usetheme{Boadilla} % Tema de layout

% Definição de cores personalizadas
\definecolor{primaryColor}{RGB}{20,45,105} % Cor primária
\definecolor{secondaryColor}{RGB}{0,100,160} % Cor secundária

% Aplicação das cores no tema
\setbeamercolor{structure}{fg=primaryColor}
\setbeamercolor{palette primary}{bg=primaryColor, fg=white}
\setbeamercolor{palette secondary}{bg=secondaryColor, fg=white}
\setbeamercolor{title}{bg=primaryColor, fg=white} % Cor do título principal

% Cores no cabeçalho
\setbeamercolor{headline}{bg=secondaryColor, fg=white}
\setbeamercolor{section in head/foot}{bg=primaryColor, fg=white}
\setbeamercolor{subsection in head/foot}{bg=secondaryColor, fg=white}

% Configurações de cores para o rodapé
\setbeamercolor{author in head/foot}{bg=primaryColor, fg=white} % Nome
\setbeamercolor{title in head/foot}{bg=secondaryColor, fg=white} % Título
\setbeamercolor{date in head/foot}{bg=primaryColor, fg=white} % Ano
\setbeamercolor{page number in head/foot}{bg=primaryColor, fg=white} % Página

% Temas internos e externos
\useinnertheme{circles} % Tema interno
\useoutertheme{miniframes} % Tema externo
\setbeamertemplate{navigation symbols}{} % Remove símbolos de navegação

%----------------------------------------------------------------------------------------
%	INFORMAÇÕES DA APRESENTAÇÃO
%----------------------------------------------------------------------------------------

\title[AS de notícia do mercado acionário]{Análise de sentimento de notícias sobre o mercado acionário brasileiro}
\author[João Squinelato]{João Guilherme Squinelato de Melo} 
\institute[PECE/Poli]{Programa de Educação Continuada da Escola Politécnica da USP \\ (PECE/Poli)} 
\date[2024]{11 / 2024} 

%----------------------------------------------------------------------------------------

\begin{document}

%----------------------------------------------------------------------------------------
%	SLIDE DE TÍTULO
%----------------------------------------------------------------------------------------

\begin{frame}
    \begin{figure}
        \includegraphics[width=0.45\linewidth]{img/logo-pece-poli.png}
    \end{figure}
    \titlepage % Exibe o slide de título
\end{frame}

%----------------------------------------------------------------------------------------
%	SLIDE DE ÍNDICE
%----------------------------------------------------------------------------------------

\begin{frame}
    \frametitle{Sumário} % Título do slide
    \tableofcontents % Exibe o índice
\end{frame}

%----------------------------------------------------------------------------------------
%	SEÇÕES DO CORPO DA APRESENTAÇÃO
%----------------------------------------------------------------------------------------

\section{Contextualização} % Seções são adicionadas para organizar sua apresentação em blocos discretos, todas as seções e subseções são automaticamente exibidas no índice como uma visão geral da apresentação, mas NÃO são exibidas como slides separados.

%------------------------------------------------

\begin{frame}
	\frametitle{Contextualização}
     
     Mercado acionário:
    \begin{itemize}
        \item Empresas de capital aberto
        \item Hipótese do Mercado Adaptativo
        \item Notícias sobre o mercado acionário
    \end{itemize}
    \bigskip
    Análise do noticiário:
    \begin{itemize}
        \item Processamento de Linguagem Natural (PLN)
        \item Análise de Sentimentos (AS)
    \end{itemize}
	
\end{frame}

%------------------------------------------------

%------------------------------------------------

\begin{frame}
	\frametitle{Contextualização}
     
    Extração de informações textuais:
    \begin{itemize}
        \item Abordagem lexical
        \item Representação vetorial
        \item Baseada em \textit{Large Language Models} (LLMs)
    \end{itemize}
    \bigskip
    Modelos de aprendizado de máquina para AS:
    \begin{enumerate}
        \item Não supervisionado
        \item Aprendizado de Máquina (AM) clássico
        \item Aprendizado Profundo (AP)
        \item Aprendizado Híbrido (AH)
        \item Baseados em LLMs
    \end{enumerate}
	
\end{frame}

%------------------------------------------------

\section{Motivação} % Seções são adicionadas para organizar sua apresentação em blocos discretos, todas as seções e subseções são automaticamente exibidas no índice como uma visão geral da apresentação, mas NÃO são exibidas como slides separados.

%------------------------------------------------

\begin{frame}
	\frametitle{Motivação}
    Desafios para o PLN de textos não ingleses:
    \begin{itemize}
        \item Dominância de estudos na língua inglesas
        \item Foco no mercado acionário estadunidense
        \item carência de bases de dados em português
    \end{itemize}
	\bigskip
    Questões a responder:
    \begin{enumerate}
        \item Modelos de LLM aplicados ao contexto de AS de notícias sobre o mercado acionário brasileiro superam abordagens clássicas?
        \item O treinamento de modelos a partir de características textuais extraídas de LLMs superam a AS obtida apenas por LLMs?
    \end{enumerate}
\end{frame}

%------------------------------------------------

\section{Problematização} % Seções são adicionadas para organizar sua apresentação em blocos discretos, todas as seções e subseções são automaticamente exibidas no índice como uma visão geral da apresentação, mas NÃO são exibidas como slides separados.

%------------------------------------------------

\begin{frame}
	\frametitle{Problematização}

    \begin{alertblock}{Problemática}
    Dificuldade de analisar o sentimento de notícias sobre o mercado acionário brasileiro, no que tange ao PLN de línguas não inglesas.
    \end{alertblock}
    \bigskip
    \begin{block}{Objetivo do Trabalho}
     Desenvolver um modelo de AH utilizando-se de informações textuais obtidas de um LLM aplicado ao contexto da AS de notícias relacionadas ao mercado acionário brasileiro.
    \end{block}

	
\end{frame}

%------------------------------------------------

\section{Trabalhos Correlatos} % Seções são adicionadas para organizar sua apresentação em blocos discretos, todas as seções e subseções são automaticamente exibidas no índice como uma visão geral da apresentação, mas NÃO são exibidas como slides separados.

%------------------------------------------------

\begin{frame}
	\frametitle{Trabalhos Correlatos}
    Santos et al. (2023):
    \begin{itemize}
        \item FinBERT-PT-BR
        \item Base de dados de notícias com polaridade anotada
        \item SentFinBERT-PT-BR
        \item Superaram o estado-da-arte para o português
    \end{itemize}
	
\end{frame}

%------------------------------------------------

\begin{frame}
	\frametitle{Trabalhos Correlatos}
    Kokab et al. (2022):
    \begin{itemize}
        \item Modelo de AH baseado em CNNs e RNNs
        \item Extração de informações textuais por meio de LLMs
        \item Superaram o estado-da-arte AS de avaliações
    \end{itemize}
	
\end{frame}

%------------------------------------------------

\section{Prováveis Contribuições} % Seções são adicionadas para organizar sua apresentação em blocos discretos, todas as seções e subseções são automaticamente exibidas no índice como uma visão geral da apresentação, mas NÃO são exibidas como slides separados.

%------------------------------------------------

\begin{frame}
	\frametitle{Prováveis Contribuições}
    No contexto do mercado acionário brasileiro, almeja-se:
    \begin{enumerate}
        \item Criar uma base de notícias anotados
        \item Contribuir para o estado-da-arte quanto à AS por meio de AH e LLMs
    \end{enumerate}
	
\end{frame}

%------------------------------------------------

\section{Referências} % Seções são adicionadas para organizar sua apresentação em blocos discretos, todas as seções e subseções são automaticamente exibidas no índice como uma visão geral da apresentação, mas NÃO são exibidas como slides separados.

\begin{frame}{Principais referências}
    \footnotesize
    \begin{thebibliography}{99}
        \bibitem{autorartigo} ASHTIANI, M. N.; RAAHEMI, B. \textit{News-based intelligent prediction of financial markets using text mining and machine learning: A systematic literature review}. Expert Systems with ApplicationsElsevier Ltd, 1 maio 2023.
        \medskip
        \bibitem{autorartigo} MAO, Y.; LIU, Q.; ZHANG, Y. \textit{Sentiment analysis methods, applications, and challenges: A systematic literature review}. Journal of King Saud University - Computer and Information SciencesKing Saud bin Abdulaziz University, , 1 abr. 2024.
        \medskip
        \bibitem{autorartigo} SOLAIMAN, S. et al. \textit{Unveiling Market Sentiments: A Comprehensive Analysis of Stock Market Responses to Diverse News Events Using Data Mining Techniques}. Proceedings of the 2024 18th International Conference on Ubiquitous Information Management and Communication, IMCOM 2024. Institute of Electrical and Electronics Engineers Inc., 2024. 
        \medskip
        \bibitem{autorartigo} SINDHU, B. et al. \textit{The Evolution of Large Language Model: Models, Applications and Challenges}. Proceedings - 2024 International Conference on Current Trends in Advanced Computing, ICCTAC 2024, 2024. 
    \end{thebibliography}
    
\end{frame}

\begin{frame}{Principais referências}
    \footnotesize
    \begin{thebibliography}{99}
        \bibitem{autorartigo} CHAHAR, R.; DUBEY, A. K. \textit{Machine learning-based Sentiment Analysis: A Comprehensive Review}. 2024 International Conference on Cognitive Robotics and Intelligent Systems, ICC - ROBINS 2024, p. 52–57, 2024.
        \medskip
        \bibitem{autorartigo} JIN, X.; CHEN, C.; YANG, X. \textit{The effect of international media news on the global stock market}. International Review of Economics and Finance, v. 89, p. 50–69, 1 jan. 2024.
        \medskip
        \bibitem{autorartigo} SANTOS, L. L.; BIANCHI, R. A. C.; COSTA, A. H. R. \textit{FinBERT-PT-BR: Análise de Sentimentos de Textos em Português do Mercado Financeiro}. Journal of the Brazilian Computer Society, 6 ago. 2023.
        \medskip
        \bibitem{autorartigo} TABINDA KOKAB, S.; ASGHAR, S.; NAZ, S. \textit{Transformer-based deep learning models for the sentiment analysis of social media data}. Array, v. 14, 1 jul. 2022. 
    \end{thebibliography}
\end{frame}


%----------------------------------------------------------------------------------------
%	SLIDE DE ENCERRAMENTO
%----------------------------------------------------------------------------------------

% \begin{frame}
%     \begin{center}
%         {\Huge Obrigado pela atenção!}
%     \end{center}
% \end{frame}

%----------------------------------------------------------------------------------------

\end{document}
